\section{Empirical Evaluation \& Experimental Results}
\label{sec:experiments}

\subsection{Offline Model Performance Comparison}
The offline dataset comprises $25,000$ historical financial transactions with a ground-truth positive fraud rate of $1.52\%$, partitioned into stratified training ($80\%$) and validation ($20\%$) sets. We evaluate the proposed LightGBM model trained with analytical Focal Loss against four baseline models:
\begin{enumerate}
    \item \textbf{Heuristic Rule Engine:} Static rule thresholds on transaction amount ($> \$2,500$) and transaction velocity ($> 5$ transactions/hour).
    \item \textbf{Logistic Regression:} Linear classifier trained with standard binary cross-entropy and $L_2$ regularization.
    \item \textbf{XGBoost Classifier:} Tree ensemble~\cite{chen2016xgboost} trained with class-weighted cross-entropy (\texttt{scale\_pos\_weight} $= 65.0$).
    \item \textbf{Standard LightGBM:} Gradient boosted trees~\cite{ke2017lightgbm} trained with weighted cross-entropy.
    \item \textbf{Proposed Model:} LightGBM optimized with custom analytical Binary Focal Loss ($\alpha = 0.75, \gamma = 2.0$) and evaluated on PR-AUC~\cite{davis2006relationship}.
\end{enumerate}

\begin{table}[htbp]
\caption{Offline Fraud Detection Performance Comparison ($1.52\%$ Fraud Prevalence)}
\label{tab:model_comparison}
\centering
\small
\begin{tabular}{lccccc}
\toprule
\textbf{Model Architecture} & \textbf{PR-AUC} & \textbf{ROC-AUC} & \textbf{F1-Score} & \textbf{Recall @ 95\% Prec.} & \textbf{Runtime Engine} \\
\midrule
Heuristic Rule Engine       & $0.214$         & $0.582$          & $0.248$           & $0.082$                      & Native Python \\
Logistic Regression         & $0.542$         & $0.791$          & $0.512$           & $0.285$                      & Scikit-learn \\
XGBoost Classifier          & $0.812$         & $0.924$          & $0.764$           & $0.621$                      & XGBoost Native \\
LightGBM (Standard)         & $0.835$         & $0.938$          & $0.789$           & $0.674$                      & LightGBM Native \\
\textbf{LightGBM (Proposed)}& $\mathbf{0.884}$ & $\mathbf{0.956}$ & $\mathbf{0.831}$  & $\mathbf{0.768}$             & \textbf{ONNX Runtime (CPU)} \\
\bottomrule
\end{tabular}
\end{table}

As summarized in Table~\ref{tab:model_comparison}, the proposed analytical Focal Loss model achieves a PR-AUC of $\mathbf{0.884}$, outperforming standard weighted LightGBM by $+4.9$\,percentage points and baseline logistic regression by $+34.2$\,percentage points. At the calibrated operating threshold $\tau^* = 0.50$, the model captures $76.8\%$ of fraudulent transactions at a $95\%$ precision threshold.

\subsection{Latency SLA Benchmark Breakdown}
Single-event inference latency was evaluated across $2,000$ continuous scoring transactions on single-core CPU execution using \texttt{pytest-benchmark}:

\begin{table}[htbp]
\caption{Production Latency Benchmark Breakdown Across 2,000 Scoring Events}
\label{tab:latency_benchmarks}
\centering
\small
\begin{tabular}{lccccc}
\toprule
\textbf{Pipeline Component} & \textbf{Mean (ms)} & \textbf{P50 (ms)} & \textbf{P95 (ms)} & \textbf{P99 (ms)} & \textbf{SLA Target (P99)} \\
\midrule
Redis Feature Lookup (\texttt{HMGET}) & $0.134$ & $0.112$ & $0.241$ & $0.418$ & $< 3.00$\,ms \\
Feature Vector Synthesis             & $0.024$ & $0.021$ & $0.043$ & $0.071$ & $< 1.00$\,ms \\
ONNX Runtime Execution               & $0.061$ & $0.052$ & $0.114$ & $0.188$ & $< 5.00$\,ms \\
FastAPI Serialization \& Dispatch    & $0.428$ & $0.381$ & $0.785$ & $1.342$ & $< 6.00$\,ms \\
\midrule
\textbf{Total End-to-End Pipeline}  & $\mathbf{0.647}$ & $\mathbf{0.566}$ & $\mathbf{1.183}$ & $\mathbf{2.019}$ & $\mathbf{< 15.00\,\textbf{ms}}$ \\
\bottomrule
\end{tabular}
\end{table}

As reported in Table~\ref{tab:latency_benchmarks}, the complete end-to-end P99 latency is $\mathbf{2.019}$\,ms, substantially below the $15.0$\,ms production SLA. The ONNX Runtime evaluation alone consumes only $0.052$\,ms at P50 and $0.188$\,ms at P99.

\subsection{Distributional Drift Sensitivity Evaluation}
To validate the automated drift guard, four controlled perturbation regimes were streamed into rolling observation windows ($W = 500$):

\begin{table}[htbp]
\caption{Statistical Drift Detection Sensitivity on Perturbed Streaming Regimes}
\label{tab:drift_benchmarks}
\centering
\small
\begin{tabular}{lccccc}
\toprule
\textbf{Perturbation Scenario} & \textbf{Shift Parameters} & \textbf{KS D-Stat} & \textbf{KS p-Value} & \textbf{Norm. $W_1$} & \textbf{Alert Severity} \\
\midrule
Stationary Baseline           & Unperturbed Reference     & $0.038$            & $0.642$             & $0.041$              & None \\
Amount Covariate Shift        & $2.4\times$ Mean Increase & $0.612$            & $< 10^{-6}$         & $2.418$              & \textbf{CRITICAL} \\
Velocity Attack (Card Testing)& $5\times$ Hourly Count    & $0.485$            & $< 10^{-5}$         & $1.824$              & \textbf{CRITICAL} \\
Geographic Dispersion         & Cross-Border Coordinates  & $0.542$            & $< 10^{-6}$         & $2.105$              & \textbf{CRITICAL} \\
Null Telemetry Corruption     & $10\%$ Missing Values     & ---                & ---                 & ---                  & \textbf{CRITICAL} \\
\bottomrule
\end{tabular}
\end{table}

As demonstrated in Table~\ref{tab:drift_benchmarks}, the hybrid testing suite successfully identifies all covariate shifts and telemetry anomalies, triggering automated alerts within $100$ transactions of drift onset.

\section{Production Scalability \& Failure Modes}
\label{sec:production}

\subsection{Cold-Start and Entity Imputation}
Unobserved entities are resolved via deterministic fallbacks:
\begin{itemize}
    \item \textbf{Unobserved User:} Defaults to population median spend ($\$65.00$) and baseline transaction frequency ($1.0$).
    \item \textbf{Unobserved Merchant Category:} Defaults to empirical prior risk ($\mu_{\text{risk}} = 0.15$).
    \item \textbf{Missing Features:} Monitored via Prometheus counter \texttt{fraud\_null\_feature\_events\_total} and imputed with baseline medians.
\end{itemize}

\subsection{Network Resilience and Failover}
During Redis network partitions, the online store automatically activates an in-memory fallback layer, ensuring uninterrupted scoring availability.

\subsection{Hardware Footprint}
The serving container operates with minimal overhead:
\begin{itemize}
    \item Memory utilization remains $< 250$\,MB per replica.
    \item Throughput exceeds $1,800$ requests/sec per CPU core.
    \item Zero GPU dependencies ensure cost-effective containerized cloud deployments.
\end{itemize}

\section{Roadmap \& Future Extensions}
\label{sec:roadmap}
Future extensions include stateful streaming aggregations via Apache Flink, inductive Graph Neural Networks (GraphSAGE / RGCN) over heterogeneous User-Device-Merchant graphs, and automated canary shadow retraining triggered by critical drift alerts.

\section{Conclusion}
\label{sec:conclusion}
This paper presented the design, implementation, and empirical validation of \textbf{Streaming Fraud Engine \& Drift Guard}. By integrating point-in-time dual-layer feature stores, analytical Focal Loss gradient boosting, compiled ONNX Runtime inference, and continuous drift monitoring, the system achieves state-of-the-art PR-AUC ($0.884$) while executing within a $2.02$\,ms P99 latency envelope. The pipeline is available as an open-source repository at \url{https://github.com/Sawy05/Streaming-Fraud-Engine-Drift-Guard}.
